% ===========================================================================
% SOURCES: docs/claude_project/08_OPEN_ITEMS.md (§5 open findings, §6 risks,
%            §7 the suggested order of work) - a snapshot of 2026-09-25: its
%            status ("never run on a board", "accuracy not recorded") is out
%            of date
%          docs/claude_project/10_TIMELINE.md (the closing observations; it
%            ends in September)
%          docs/design_decisions.md: DD-034, DD-035, DD-036, and
%            docs/lab_notebook.md from 2026-10-07 on, for everything later
% ===========================================================================
\chapter{Conclusion and Future Work}
\label{chap:conclusion}

\section{Conclusion}
\label{sec:conclusion}
% Answer Section~\ref{sec:objectives} against the tiers of \DD{004}, in order, without
% hedging:
%   Foundation - REACHED. Full MobileNetV2 bit-exact in simulation,
%     \thElements elements.
%   Core - REACHED on 2026-10-07: the bare-metal application on the ZCU104
%     reproduced the known answer bit for bit (Section~\ref{sec:board-remaining}).
%     Say exactly what that run covered - one image - and what the other
%     bring-up checks showed, once they are recorded.
%   Stretch - built and tested on the PC (\DD{035}). Write the board result
%     here once it exists (Section~\ref{sec:camera-result}), or one clause
%     saying it did not run.
%
% Then the comparison (\DD{036}), answered as plainly as the tiers: against a
% Raspberry Pi 4 and a Jetson Nano running the same model on the same images,
% where the accelerator is ahead, where it is not, and by how much - latency
% and energy per inference. Do not soften a loss: a measured loss on one axis
% is a result too, and the committee will trust the wins more for it.
%
% Then the three observations the lab notebook supports, which are the
% methodological conclusions and the most transferable part of the thesis:
%   1. The software model was finished before any RTL was written. It changed
%      exactly once afterwards, by decision - the leaf-split retrain
%      (\DD{034}), redeployed to every platform (\DD{036}) - and the hardware
%      did not have to change: the weights are a DDR blob, the accumulator
%      bound was rechecked by the exporter rather than remembered
%      (Section~\ref{sec:acc-bound}), and the golden vectors re-proved the
%      whole chain. That is what made every later step a comparison rather
%      than a negotiation.
%      TODO(REDEPLOY): this wording is true once the redeploy is done and the
%      simulations pass unchanged; until then "never changed afterwards" is
%      the accurate sentence.
%   2. Each verification level found bugs the level below structurally could
%      not, and each was added BEFORE the phase that needed it.
%   3. The estimates were mostly right and the assumptions were not. The cycle
%      model came within 7%; the DSP arithmetic was right. What was wrong was
%      assuming int8 packing existed, assuming the tool would choose URAM,
%      assuming a measured accumulator bound would survive a retrain, and
%      assuming a shared stage could be costed against its cheapest consumer.
%      EVERY SIGNIFICANT SURPRISE IN THIS PROJECT WAS AN ASSUMPTION, NOT AN
%      ARITHMETIC ERROR. That is the sentence to end on.

\section{Future Work}
\label{sec:future}
% Concrete and sized, not a wish list. Each item should say what it costs and
% what it buys:
%
%   - (The board run used to head this list. It happened on 2026-10-07 -
%     Section~\ref{sec:board-remaining}. Only bring-up checks still unrecorded
%     at submission belong here now.)
%   - \DD{023}, the act_buffer read pipeline: the stated route to 250 MHz, but
%     Section~\ref{sec:board-impl} found the in-context critical path is elsewhere -
%     replicating wgt_dma's 216-bit output register per destination buffer is
%     cheaper and comes first.
%   - C-11, weight double-buffering. Report the corrected figure honestly: the
%     65% overhead was computed against a fill rate nobody had measured, and
%     the real DMA gives \thLoadOverheadPct\% network-wide. 65% still stands for
%     features.18.0 alone. Double-buffering can hide at most \thLoadOverheadPct\%.
%   - The remaining open review findings: C-10, T-6, T-8, T-9
%     (Section~\ref{sec:rtl-review}).
%   - Generality. The instruction ceiling is \thInstructions/\thInstructions{} with
%     zero spare, and the weight address width has 22% headroom FOR THIS MODEL -
%     features.18.0 is 800/1024 and the classifier is small only because there
%     are \thClasses{} classes. A 1000-class head would need 2,560.
%   - The Stretch tier is built (\DD{035}); whatever of it has not run on the
%     board by submission goes here, sized.
%   - Field data: out of scope by decision (\DD{036}). Photographs from the
%     grapevine IoT deployment [morellos2024iot] - first as a test set of their
%     own, since every accuracy figure here is on PlantVillage photographs taken
%     under controlled conditions; DD-034 names them as the thing that reopens
%     the accuracy result. Retraining on them comes after that.
%   - An int8 GPU row: a Jetson with int8 tensor cores (an Orin) would give the
%     comparison the column the Nano cannot (\DD{036}).
%   - A cycle counter in the CSR, so the board reports its own latency instead
%     of the PS timing it (Section~\ref{sec:throughput}).
%   - The vendor DPU, measured on the same board - if \DD{036} leaves it as a
%     discussion point.
